\section{分布式训练：五种并行策略}

\subsection{并行化的四个维度}

Transformer 模型本质上是 $L$ 层的堆叠，每层操作一个三维张量 $(\text{batch}, \text{sequence}, \text{dim})$。这为我们提供了四个天然的并行化轴：

\begin{figure}[H]
\centering
\includegraphics[width=0.95\textwidth]{slides-images/slide-012.jpg}
\caption{Transformer 的四个并行化维度：Data（batch）、Context（sequence）、Tensor（dim）、Pipeline（layers）\protect\footnotemark}
\end{figure}
\footnotetext{Slides 第12页。}

\begin{importantbox}{五种并行策略}
\begin{enumerate}
    \item \textbf{数据并行（Data Parallelism, DP）}：在 batch 维度上切分
    \item \textbf{全分片数据并行（FSDP）}：DP + 模型权重分片
    \item \textbf{混合分片数据并行（HSDP）}：FSDP + DP 的二维组合
    \item \textbf{上下文并行（Context Parallelism, CP）}：在 sequence 维度上切分
    \item \textbf{流水线并行（Pipeline Parallelism, PP）}：在 layers 维度上切分
    \item \textbf{张量并行（Tensor Parallelism, TP）}：在 dim 维度上切分
\end{enumerate}
\end{importantbox}

\subsection{本章小结}

Transformer 模型的四个维度（batch、sequence、dim、layers）为分布式训练提供了天然的并行化轴。不同维度的并行策略具有不同的通信需求和适用场景，需要根据模型规模、GPU 数量和集群拓扑进行选择。

%% ========== 数据并行 ==========

